\documentclass[12pt,a4paper]{article}

% ================= PACKAGES =================
\usepackage[margin=1in]{geometry}
\usepackage[T1]{fontenc}
\usepackage{graphicx}
\usepackage{titlesec}
\usepackage{tabularx}
\usepackage{longtable}
\usepackage{array}
\usepackage{hyperref}
\usepackage{listings}
\usepackage{xcolor}
\usepackage{enumitem}
\usepackage{setspace}
\usepackage[most]{tcolorbox}
\usepackage{mathptmx}   % Times New Roman-equivalent font

% ================= FONT SIZE LADDER =================
\newcommand{\sizeEighteen}{\fontsize{18pt}{22pt}\selectfont}  % \section
\newcommand{\sizeSixteen}{\fontsize{16pt}{20pt}\selectfont}   % \subsection
\newcommand{\sizeFourteen}{\fontsize{14pt}{17pt}\selectfont}  % \subsubsection
\newcommand{\sizeTwelve}{\fontsize{12pt}{14.5pt}\selectfont}  % body / minimum

% ================= HEADER COLORS =================
\definecolor{RamaiahBlue}{RGB}{27,31,92}
\definecolor{RamaiahMaroon}{RGB}{139,30,63}

% ================= CODE LISTING STYLE =================
\lstdefinestyle{codestyle}{
    backgroundcolor=\color{black!5},
    commentstyle=\color{green!50!black},
    keywordstyle=\color{blue},
    stringstyle=\color{orange},
    basicstyle=\ttfamily\fontsize{10pt}{12.5pt}\selectfont,
    numberstyle=\ttfamily\fontsize{9pt}{11pt}\selectfont,
    breaklines=true,
    breakatwhitespace=false,
    frame=single,
    numbers=left,
    numbersep=8pt,
    xleftmargin=2.5em,
    aboveskip=8pt,
    belowskip=8pt,
    showstringspaces=false
}
\lstset{style=codestyle}

% ================= HEADING FORMATS =================
\titleformat{\section}{\sizeEighteen\bfseries}{\thesection}{1em}{}
\titleformat{\subsection}{\sizeSixteen\bfseries}{\thesubsection}{1em}{}
\titleformat{\subsubsection}{\sizeFourteen\bfseries}{\thesubsubsection}{1em}{}

\titlespacing*{\section}{0pt}{18pt}{10pt}
\titlespacing*{\subsection}{0pt}{16pt}{8pt}
\titlespacing*{\subsubsection}{0pt}{14pt}{7pt}

% ================= REPEATABLE HEADER BANNER =================
\newcommand{\reportheader}{%
\noindent
\begin{minipage}[c]{0.42\textwidth}
  \IfFileExists{ramaiah_logo.jpg}{\includegraphics[height=1.2cm]{ramaiah_logo.jpg}}{\fbox{\rule{0pt}{1.2cm}\textbf{ Ramaiah Logo }}}
\end{minipage}%
\begin{minipage}[c]{0.58\textwidth}
  \raggedleft
  {\color{RamaiahBlue}\sizeFourteen\bfseries Department of CSE (Artificial Intelligence \& Machine Learning) and CSE (Cyber Security)}
\end{minipage}\\[0.2cm]
\noindent\textcolor{RamaiahMaroon}{\rule{\linewidth}{1.2pt}}
\vspace{0.4cm}
}

\begin{document}
\pagestyle{empty}
\sizeTwelve

% ======================================================
% PAGE 1: TITLE PAGE
% ======================================================
\reportheader

\vspace{0.8cm}
\begin{center}
{\sizeEighteen \bfseries SELF LEARNING ACTIVITY REPORT \par}

\vspace{1.0cm}

{\sizeTwelve Submitted in partial fulfilment of the Requirements for \\[0.2cm]
the Academics of 5\textsuperscript{th} Semester, B.E in CSE (Cyber Security) \par}

\vspace{1.0cm}

{\sizeEighteen \bfseries DEVOPS \par}

\vspace{0.4cm}

{\sizeSixteen \bfseries 24CYAEC59 \par}

\vspace{1.0cm}

{\sizeTwelve In \par}

\vspace{0.35cm}

{\sizeFourteen \bfseries Dept. of CSE (CYBER SECURITY) \par}

\vspace{1.0cm}

{\sizeTwelve By \par}

\vspace{0.35cm}

{\sizeFourteen \bfseries Kavya Poddar \par}
\vspace{0.15cm}
{\sizeFourteen \bfseries USN: 1MS24CY\_\_\_ \par}

\vspace{1.2cm}

{\sizeTwelve Under the guidance of \par}
\vspace{0.25cm}

{\sizeFourteen \bfseries Mr. Prashanth Reddy C \par}
\vspace{0.15cm}

{\sizeTwelve Senior Lecturer \par}
{\sizeTwelve Department of CSE (AI \& ML) and CSE (Cyber Security) \par}

\vspace{1.2cm}

{\sizeSixteen \bfseries M. S. Ramaiah Institute of Technology \par}
\vspace{0.25cm}
{\sizeTwelve (Autonomous Institute Affiliated to VTU) \\ MSR Nagar, MSRIT Post \par}
\vspace{0.1cm}
{\sizeTwelve Bangalore-560054, Karnataka, India \par}
\end{center}

\newpage

% ======================================================
% PAGE 2: CERTIFICATE
% ======================================================
\reportheader

\vspace{1.0cm}
\begin{center}
{\sizeEighteen \bfseries Certificate \par}
\end{center}

\vspace{0.8cm}

\noindent
{\sizeTwelve
This is to certify that the work carried out is a bonafide record of the project based learning activity work on the DevOps course done by \textbf{Kavya Poddar (USN: 1MS24CY\_\_\_)} under my supervision and guidance, in partial fulfilment of the requirements for the award of the Bachelor of Engineering in Computer Science \& Engineering (Cyber Security) from Ramaiah Institute of Technology, Bangalore for the academic year 2026-2027.
\par}

\vspace{3.5cm}

\noindent
{\sizeTwelve Candidate Sign \hfill Faculty Sign \par}

\newpage
\pagestyle{plain}

% ======================================================
% LIST OF ACTIVITIES
% ======================================================
\section*{LIST OF PROJECT BASED LEARNING COMPONENTS}

\noindent
{\sizeTwelve
\textbf{COURSE NAME:} DevOps \hfill \textbf{COURSE CODE:} 24CYAEC59 \\
\textbf{CREDITS:} 1:0:0 \hfill \textbf{CONTACT HOURS:} 15L \\
\textbf{MAX. MARKS:} 20
\par}

\vspace{0.6cm}

{\sizeTwelve
\setlength{\tabcolsep}{4pt}
\renewcommand{\arraystretch}{1.25}
\begin{longtable}{|>{\centering\arraybackslash}m{1.2cm}|>{\raggedright\arraybackslash}p{10.6cm}|>{\centering\arraybackslash}m{1.8cm}|}
\hline
\multicolumn{1}{|c|}{\textbf{Sl. No}} &
\multicolumn{1}{c|}{\textbf{Description}} &
\multicolumn{1}{c|}{\textbf{Marks}} \\
\hline
\endhead

1 &
\textbf{CI/CD Pipeline using Git, Python/Pytest and Jenkins.}
Set up a version-controlled build pipeline for a sample application. Create a Git repository, use feature branching and merging, configure dependency management and testing, integrate the project into a Jenkins job that pulls from SCM, triggers automatically on commit, executes linting, runs automated tests with coverage, performs build verification, and adds multi-stage quality gate enhancements.

\textit{Covers: Unit 1 (DevOps culture/pipeline), Unit 2 (Git, Dependency Management), Unit 3 (Jenkins CI). Maps to CO1, CO2, CO3.}
& 10 \\
\hline

2 &
\textbf{Configuration Management with Ansible and Deployment on Kubernetes.}
Write an Ansible playbook (using variables, handlers, and the \texttt{app\_setup} role) to configure and deploy the application built in Component 1. Containerize the application using Docker and write Kubernetes manifests (Deployment and NodePort Service). Deploy on a local cluster (Minikube), demonstrate scaling using replica sets, show the use of labels and selectors, and demonstrate deployment self-healing via Liveness and Readiness probes.

\textit{Covers: Unit 4 (Ansible), Unit 5 (Kubernetes). Maps to CO4, CO5.}
& 10 \\
\hline

\end{longtable}
\par}

\newpage

% ======================================================
% ASSESSMENT RUBRICS
% ======================================================
\section*{Assessment Rubrics}
{\sizeTwelve Evaluation for \textbf{each} 10-mark component will be based on the following rubrics.\par}

\vspace{0.4cm}

{\sizeTwelve
\setlength{\tabcolsep}{4pt}
\renewcommand{\arraystretch}{1.25}
\begin{longtable}{|>{\raggedright\arraybackslash}p{3.0cm}|>{\raggedright\arraybackslash}p{3.5cm}|>{\raggedright\arraybackslash}p{3.5cm}|>{\raggedright\arraybackslash}p{3.5cm}|}
\hline

\textbf{Criteria} & \textbf{2 --- Good} & \textbf{1 --- Average} & \textbf{0 --- Poor} \\

\hline
\endhead
Concept Understanding (2 marks) & Answers ``why'' questions correctly and immediately (e.g., webhook vs.\ polling, what a handler does, Deployment vs.\ ReplicaSet, why health probes matter) & Gets the basic idea right but stumbles on follow-up questions or confuses related concepts & Cannot explain own pipeline; answers suggest it was copied without understanding \\
\hline
Core Implementation (2 marks) & Required pipeline (Git+Jenkins for Comp.\ 1 / Ansible+K8s+scaling for Comp.\ 2) runs end-to-end live, without intervention & Works but needs manual triggering, a workaround, or partial failure & Doesn't run at all during the demo; only screenshots are shown \\
\hline
Enhancement (2 marks) & At least one enhancement implemented and demonstrated live, working correctly & Enhancement attempted but incomplete, or only explained without being runnable & No enhancement attempted \\
\hline
Analysis \& Interpretation (2 marks) & Explains build logs, scaling behaviour, or enhancement outcomes in own words; interprets output unprompted & Describes what happened but not why, or needs prompting to interpret output & No interpretation offered; treats commands as a checklist \\
\hline
Documentation \& Presentation (2 marks) & Report fully personalized and confidently walks through the pipeline unprompted & Mostly complete report or needs guidance/prompting during the demo & Unedited template, mismatched screenshots, or cannot navigate own setup \\
\hline
\end{longtable}
\par}

\newpage

% ======================================================
% MARKS AWARDED TABLE
% ======================================================
\section*{Marks Awarded}

\vspace{0.3cm}

{\sizeTwelve
\setlength{\tabcolsep}{3pt}
\renewcommand{\arraystretch}{1.5}
\begin{longtable}{|>{\centering\arraybackslash}p{1.0cm}|>{\raggedright\arraybackslash}p{2.2cm}|>{\raggedright\arraybackslash}p{2.2cm}|>{\raggedright\arraybackslash}p{2.0cm}|>{\raggedright\arraybackslash}p{2.3cm}|>{\raggedright\arraybackslash}p{2.5cm}|>{\centering\arraybackslash}p{1.4cm}|}
\hline
\textbf{Sl. No} &
\textbf{Concept Under- standing (2)} &
\textbf{Core Imple- mentation (2)} &
\textbf{Enhance- ment (2)} &
\textbf{Analysis \& Interpre- tation (2)} &
\textbf{Documen- tation \& Presenta- tion (2)} &
\textbf{Total} \\
\hline
\parbox[c][1.1cm][c]{1.0cm}{\centering 1} & & & & & & \\
\hline
\parbox[c][1.1cm][c]{1.0cm}{\centering 2} & & & & & & \\
\hline
\multicolumn{6}{|c|}{\parbox[c][1.1cm][c]{13.0cm}{\centering\textbf{TOTAL MARKS}}} &
\parbox[c][1.1cm][c]{1.4cm}{\centering} \\
\hline
\end{longtable}
\par}

\newpage

% ======================================================
% COMPONENT 1 FULL REPORT
% ======================================================
\section*{Component 1: CI/CD Pipeline using Git, Python/Pytest and Jenkins}

\subsection*{Abstract}
This activity implements an automated Continuous Integration (CI) pipeline using Git, Python 3.12, Pytest, and Jenkins for a high-performance FastAPI microservice (\texttt{todo-api}). Source code modifications are tracked through feature branch workflows and merged into \texttt{main}. The Jenkins declarative pipeline orchestrates automated code checkout, isolated virtual environment creation, code quality linting with Flake8, execution of a 14-test suite with code coverage, and build verification upon every detected code commit.

\subsection*{Introduction}
DevOps bridges the gap between software development and IT operations through automated validation mechanisms. In this implementation, a RESTful TODO API was constructed using the FastAPI ASGI framework with Pydantic data schemas. Git provides distributed version control allowing isolated feature development. Jenkins monitors the remote repository and executes a multi-stage validation pipeline on every push. This rapid feedback loop guarantees that regressions, syntax violations, and integration defects are detected immediately before software is packaged.

\subsection*{Implementation}

\subsubsection*{Development Environment}
\begin{itemize}[nosep]
    \item \textbf{Operating System:} macOS / Linux
    \item \textbf{Version Control:} Git 2.50.x \& GitHub
    \item \textbf{Programming Language \& Framework:} Python 3.12, FastAPI, Uvicorn
    \item \textbf{Testing \& Linting Frameworks:} Pytest 8.x, HTTPX TestClient, Pytest-Cov, Flake8
    \item \textbf{CI Automation Server:} Jenkins 2.x LTS (Docker container)
    \item \textbf{Application:} \texttt{todo-api} (REST microservice with in-memory CRUD \& health checks)
\end{itemize}

\subsubsection*{Step 1: Git Repository and Branching Workflow}
The repository was initialized under Git version control. Features were isolated on branches before being merged into the production branch \texttt{main} and pushed upstream to GitHub.

\begin{lstlisting}[language=bash]
# Initialize repository and create feature branch
cd devops-demo
git checkout -b feature/todo-fastapi-service

# Commit application source, dependencies and tests
git add app/ tests/ requirements.txt pyproject.toml .gitignore README.md
git commit -m "feat: migrate backend to FastAPI REST API with pytest test suite"

# Configure CI pipeline and push to main
git add Jenkinsfile .github/
git commit -m "ci: configure Jenkins pipeline with lint, test coverage, and build verification"
git checkout main
git merge feature/todo-fastapi-service
git push origin main
\end{lstlisting}

\subsubsection*{Step 2: Dependency Management (requirements.txt)}
Project dependencies are managed declaratively using Python packaging standards.

\begin{lstlisting}[language=bash]
fastapi>=0.111.0
uvicorn[standard]>=0.29.0
pydantic>=2.0.0
pytest>=8.0.0
httpx>=0.27.0
pytest-cov>=5.0.0
\end{lstlisting}

\subsubsection*{Step 3: Jenkins Declarative Pipeline (Jenkinsfile)}
The Jenkinsfile defines five sequential stages to rigorously validate code health across clean workspace environments.

\begin{lstlisting}[language=groovy]
pipeline {
    agent any
    environment {
        VENV_DIR = '.venv'
    }
    stages {
        stage('Checkout') {
            steps {
                echo 'Checking out source code from GitHub...'
                checkout scm
            }
        }
        stage('Setup Python Environment') {
            steps {
                sh '''
                    python3 -m venv $VENV_DIR
                    . $VENV_DIR/bin/activate
                    pip install --upgrade pip --quiet
                    pip install -r requirements.txt
                '''
            }
        }
        stage('Lint') {
            steps {
                sh '''
                    . $VENV_DIR/bin/activate
                    pip install flake8 --quiet
                    flake8 app/ --max-line-length=100 --ignore=E501,W503
                '''
            }
        }
        stage('Test') {
            steps {
                sh '''
                    . $VENV_DIR/bin/activate
                    pytest tests/ -v --tb=short --cov=app --cov-report=term
                '''
            }
        }
        stage('Build Verification') {
            steps {
                sh '''
                    . $VENV_DIR/bin/activate
                    python3 -c "from app.main import app; print('App import: OK')"
                '''
            }
        }
    }
    post {
        always {
            cleanWs()
        }
    }
}
\end{lstlisting}

\subsubsection*{Step 4: CI Enhancement --- Multi-Stage Quality Gates in CI}
The selected CI enhancement is \textbf{Automated Multi-Stage Quality Gates (PEP8 Linting + Code Coverage Testing)}. By executing Flake8 before testing, code smell and style violations halt the build early. The test suite contains 14 distinct test cases verifying endpoint responses, schema validations, error handling, and coverage metrics.

\begin{lstlisting}[language=Python]
# Sample excerpt from tests/test_todos.py
def test_health_check():
    response = client.get("/health")
    assert response.status_code == 200
    assert response.json() == {"status": "ok", "service": "todo-api"}

def test_create_todo_basic():
    payload = {"title": "Buy groceries"}
    response = client.post("/todos/", json=payload)
    assert response.status_code == 201
    assert response.json()["id"] == 1
    assert response.json()["completed"] is False
\end{lstlisting}

\subsubsection*{Description of the Activity}
The pipeline was connected to the GitHub repository via Jenkins Git SCM. When changes are pushed, Jenkins initializes a clean workspace, spins up an isolated virtual environment, installs dependencies, validates code compliance with Flake8, runs all 14 unit test assertions via Pytest with coverage analysis, and performs a smoke import test. If any step fails, the pipeline aborts immediately, safeguarding production branches from defective code.

\subsection*{Results}

\begin{figure}[htbp]
\centering
% \includegraphics[width=0.88\linewidth]{screenshots/comp1_git_history.png}
\caption{Git Feature Branching and Merge Commit History (\texttt{git log --graph --oneline})}
\end{figure}

\begin{figure}[htbp]
\centering
% \includegraphics[width=0.88\linewidth]{screenshots/comp1_jenkins_pipeline.png}
\caption{Jenkins Declarative Pipeline Execution showing 5 successful green stages}
\end{figure}

\begin{figure}[htbp]
\centering
% \includegraphics[width=0.88\linewidth]{screenshots/comp1_test_results.png}
\caption{Pytest Unit Test Suite Output showing 14 passed test assertions and coverage}
\end{figure}

\subsection*{Conclusion}
Component 1 successfully automated source code integration. Enforcing automated testing and linting gates within Jenkins eliminated manual verification errors, guaranteeing rapid feedback and high code quality.

\newpage

% ======================================================
% COMPONENT 2 FULL REPORT
% ======================================================
\section*{Component 2: Configuration Management with Ansible and Deployment on Kubernetes}

\subsection*{Abstract}
This activity automates server provisioning using Ansible and orchestrates the containerized \texttt{todo-api} application on Kubernetes. An idempotent Ansible role (\texttt{app\_setup}) installs runtime packages, provisions virtual environments, and manages the daemon process with change-triggered handlers. The application is containerized with Docker and deployed to a Minikube Kubernetes cluster with 3 replicas, scaled to 5 replicas, and configured with Liveness and Readiness probes on the \texttt{/health} endpoint.

\subsection*{Introduction}
Modern cloud operations require infrastructure to be treated as code. Ansible automates host configurations without requiring client agents, using declarative playbooks and reusable roles. Kubernetes orchestrates containerized workloads across clusters, providing service discovery, horizontal pod autoscaling, and self-healing. This component demonstrates how Ansible standardizes the runtime environment while Kubernetes ensures high availability and resilience for the containerized application.

\subsection*{Implementation}

\subsubsection*{Development Environment}
\begin{itemize}[nosep]
    \item \textbf{Configuration Management:} Ansible 2.16+
    \item \textbf{Container Engine:} Docker 29.x
    \item \textbf{Local Kubernetes Cluster:} Minikube v1.32+
    \item \textbf{Cluster CLI:} kubectl v1.28+
\end{itemize}

\subsubsection*{Step 1: Ansible Playbook and Role (\texttt{app\_setup})}
The Ansible playbook targets localhost, incorporates host variables (\texttt{app\_port: 8000}, \texttt{app\_dir: /opt/todo-api}), executes modular tasks via the \texttt{app\_setup} role, and triggers handlers only upon detected changes.

\begin{lstlisting}[language=yaml]
---
# ansible/playbook.yml
- name: Setup TODO API Server
  hosts: localhost
  connection: local
  become: true
  vars:
    app_port: 8000
    app_dir: "/opt/todo-api"
  roles:
    - app_setup
\end{lstlisting}

\begin{lstlisting}[language=yaml]
# ansible/roles/app_setup/tasks/main.yml
---
- name: Ensure Python3, pip, and venv are installed
  apt:
    name: [python3, python3-pip, python3-venv]
    state: present
    update_cache: yes

- name: Ensure application directory exists
  file:
    path: "{{ app_dir }}"
    state: directory
    mode: '0755'

- name: Install Python dependencies into virtual environment
  pip:
    requirements: "{{ app_dir }}/requirements.txt"
    virtualenv: "{{ app_dir }}/venv"

- name: Copy application code
  copy:
    src: "{{ playbook_dir }}/../app"
    dest: "{{ app_dir }}/"
    mode: '0755'
  notify: Restart app
\end{lstlisting}

\begin{lstlisting}[language=yaml]
# ansible/roles/app_setup/handlers/main.yml
---
- name: Restart app
  shell: |
    if [ -f {{ app_dir }}/app.pid ]; then
      kill -9 $(cat {{ app_dir }}/app.pid) 2>/dev/null || true
    fi
    nohup {{ app_dir }}/venv/bin/uvicorn app.main:app --host 0.0.0.0 --port {{ app_port }} > {{ app_dir }}/app.log 2>&1 &
    echo $! > {{ app_dir }}/app.pid
  args:
    chdir: "{{ app_dir }}"
\end{lstlisting}

\subsubsection*{Step 2: Dockerfile Containerization}
The application is packaged into a slim, secure container image based on \texttt{python:3.12-slim}.

\begin{lstlisting}[language=bash]
FROM python:3.12-slim
WORKDIR /app
COPY requirements.txt .
RUN pip install --no-cache-dir -r requirements.txt
COPY app/ ./app/
EXPOSE 8000
CMD ["uvicorn", "app.main:app", "--host", "0.0.0.0", "--port", "8000"]
\end{lstlisting}

\subsubsection*{Step 3: Kubernetes Deployment and NodePort Service Manifests}
The manifests declare a Deployment managing 3 Pod replicas with Liveness/Readiness probes and a NodePort Service exposing port 80 externally via port 30080.

\begin{lstlisting}[language=yaml]
# k8s/deployment.yml
apiVersion: apps/v1
kind: Deployment
metadata:
  name: todo-api
  labels:
    app: todo-api
spec:
  replicas: 3
  selector:
    matchLabels:
      app: todo-api
  template:
    metadata:
      labels:
        app: todo-api
    spec:
      containers:
        - name: todo-api
          image: todo-api:latest
          imagePullPolicy: IfNotPresent
          ports:
            - containerPort: 8000
          livenessProbe:
            httpGet:
              path: /health
              port: 8000
            initialDelaySeconds: 10
            periodSeconds: 15
          readinessProbe:
            httpGet:
              path: /health
              port: 8000
            initialDelaySeconds: 5
            periodSeconds: 10
---
# k8s/service.yml
apiVersion: v1
kind: Service
metadata:
  name: todo-api-service
spec:
  selector:
    app: todo-api
  type: NodePort
  ports:
    - protocol: TCP
      port: 80
      targetPort: 8000
      nodePort: 30080
\end{lstlisting}

\subsubsection*{Step 4: Deployment and Scaling Execution}
The Kubernetes resources are applied to Minikube and dynamically scaled from 3 to 5 replicas.

\begin{lstlisting}[language=bash]
# Build image inside Minikube
eval $(minikube docker-env)
docker build -t todo-api:latest .

# Deploy manifests
kubectl apply -f k8s/deployment.yml
kubectl apply -f k8s/service.yml

# Verify initial 3 replicas
kubectl get pods -l app=todo-api

# Scale horizontally to 5 replicas
kubectl scale deployment todo-api --replicas=5

# Verify scaled pods
kubectl get pods -l app=todo-api -o wide
\end{lstlisting}

\subsubsection*{Step 5: Deployment Enhancement --- Liveness and Readiness Probes}
The selected deployment enhancement is **Liveness and Readiness Health Probes**. The readiness probe ensures traffic is routed to the Pod only after Uvicorn starts. The liveness probe continuously queries `/health`; if the service hangs, the Kubelet automatically restarts the container, achieving automated self-healing.

\subsubsection*{Description of the Activity}
The Ansible playbook configured server dependencies and handled idempotency by only firing the restart handler when the application code changed. In Kubernetes, the Deployment controller established a ReplicaSet of 3 pods. Scaling to 5 replicas caused the cluster scheduler to instantly bind and initialize two new pods. The NodePort service load-balanced incoming requests across all available endpoints.

\subsection*{Results}

\begin{figure}[htbp]
\centering
% \includegraphics[width=0.88\linewidth]{screenshots/comp2_ansible_run.png}
\caption{Ansible Playbook Run Output showing idempotent tasks and handler execution}
\end{figure}

\begin{figure}[htbp]
\centering
% \includegraphics[width=0.88\linewidth]{screenshots/comp2_k8s_scaling.png}
\caption{Kubernetes Pods Scaling from 3 to 5 Replicas (\texttt{kubectl get pods})}
\end{figure}

\begin{figure}[htbp]
\centering
% \includegraphics[width=0.88\linewidth]{screenshots/comp2_probe_describe.png}
\caption{Kubernetes Pod Description showing Liveness and Readiness Probes (\texttt{kubectl describe pod})}
\end{figure}

\subsection*{Conclusion}
Component 2 demonstrated continuous deployment and orchestration. Ansible ensured reproducible host configurations, while Kubernetes provided automated healing, zero-downtime scaling, and resilient service routing.

\newpage

% ======================================================
% REFERENCES
% ======================================================
\section*{References}

\begin{enumerate}[nosep]
    \item Jennifer Davis and Katherine Daniels, \textit{Effective DevOps}, 1st Edition, Shroff/O'Reilly Publications, 2021.
    \item Sanjeev Sharma, \textit{The DevOps Adoption Playbook}, 1st Edition.
    \item Official Git Documentation -- \url{https://git-scm.com/doc}
    \item Official Jenkins Documentation -- \url{https://www.jenkins.io/doc/}
    \item Official Ansible Documentation -- \url{https://docs.ansible.com/}
    \item Official Kubernetes Documentation -- \url{https://kubernetes.io/docs/}
    \item FastAPI Official Documentation -- \url{https://fastapi.tiangolo.com/}
\end{enumerate}

\end{document}
